\section{Related Work}
\label{sec:related}

\paragraph{Predicting behavioral and neural responses.}
A growing body of work asks whether LLM representations line up with human brain activity and behavior during language processing. Model-derived features, most often next-word surprisal or hidden-state embeddings, predict fMRI and MEG responses during naturalistic listening and reading \citep{schrimpf2021neural,caucheteux2022brains,caucheteux2023evidence}, explain variance in fMRI and ECoG recordings beyond simpler linguistic features \citep{toneva2019interpreting}, and track a hierarchy of predictive signals across the cortical language network \citep{caucheteux2023evidence}. These correspondences hold for models trained on child-scale corpora \citep{hosseini2024artificial}, can be strengthened by instruction tuning \citep{aw2024instructiontuning}, and are matched or exceeded by recurrent architectures on several comprehension metrics \citep{michaelov2024revenge}. Alignment does not grow indefinitely with training, however. Tracking eight Pythia models (14M to 6.7B parameters) across 34 training checkpoints on five neural benchmarks, \citet{alkhamissi-etal-2025-language} find that alignment with the human language network rises sharply early in training, peaks after roughly 8B tokens and then saturates, following formal linguistic competence (knowledge of linguistic rules) more closely than functional competence, which keeps improving. They also report that, once controlling for the number of features, model size is not a reliable predictor of brain alignment, and that the correlation of brain alignment with next-word prediction and with behavioral alignment fades once models surpass human language proficiency. Intermediate layers already encode next-token distributions in a way that resembles human online processing \citep{kuribayashi2025humanlike}, and reading times can be decoded from hidden states of models never trained to predict them \citep{tsipidi2026probing}. Taken together, these findings motivate the Platonic Representation Hypothesis \citep{huh2024platonic}: models trained with different objectives, architectures, and data settle into overlapping representational geometry because they solve related tasks over similar statistics.

\paragraph{Surprisal and reading behavior.}
Surprisal links language-model predictions to processing difficulty \citep{hale2001probabilistic,levy2008expectation,smithlevy2013effect} and remains the dominant predictor of reading times \citep{goodkind2018predictive,wilcox2020predictive,wilcox2023testing}, with psychometric predictive power now benchmarked across model families and sizes \citep{kuribayashi2024psychometric}. Its relationship with language-model quality is not monotonic: several studies report that smaller or less accurate models fit reading times better than large ones \citep{ohschuler2023why,ohschuler2023twobillion,kuribayashi2021lower,wilcox2024bigger}. Related work has generalized what counts as an anticipation signal in online processing \citep{giulianelli2024generalized,opedal2024role}. Neural encoding studies, by contrast, typically find that larger models predict brain responses better \citep{schrimpf2021neural,caucheteux2022brains}. These evaluations need not favor the same model, which is why our design keeps behavioral and neural outcomes separate and evaluates both embeddings and surprisal within each domain.

\paragraph{Memory and developmental constraints.}
If shared task pressure suffices to produce convergence, why constrain a model at all? One answer, pursued mostly in psycholinguistics, is that human sentence processing operates under resource limits that standard LLMs do not share, and that building those limits in should make a model behave more like a human reader. On the architecture side, this has taken the form of restricting attention itself: single-head or otherwise capacity-limited attention \citep{timkey2023limited}, distance-based decay that biases attention toward nearby tokens \citep{press2021train,devarda2024locally,clark2025linear,thamma2025humanlike}, content-sensitive forget gates \citep{lin2025forgetting} or stick-breaking attention \citep{tan2025scaling}, and memory-bounded or resource-rational formulations of the encoding process \citep{hahn2022resourcerational,lieder2019resourcerational,xu2026memory,xufutrell2026strategic,rasmussen2017leftcorner,keshev2025working}. These choices are usually motivated by capacity-limited theories of sentence processing, such as activation-based memory retrieval \citep{lewisvasishth2005activation} and lossy-context surprisal \citep{futrell2020lossy}. On the data side, the BabyLM Challenge leaves the architecture alone but restricts training exposure to the 10--100 million words a child hears by early adolescence \citep{warstadt2023cfp,warstadt2023findings,hu2024secondbabylm}, building on earlier work on pretraining with child-directed language \citep{huebner2021babyberta}. We organize these approaches as two factors and explicitly estimate their interaction. Because the evaluated models also differ in architecture and training details, this grouping supports comparative associations rather than an isolated causal test of either constraint.

\paragraph{Prediction and representational correspondence.}
Encoding models assess information recoverable through a learned mapping. RSA instead compares relations among stimuli in model and response spaces \citep{kriegeskorte2008rsa}. These are distinct tests: a representation may support accurate decoding without matching response geometry under a specified distance metric, and a learned mapping has a capacity that grows with the width of the representation. Feature scaling also affects embedding distances \citep{timkey2021allbark}. We therefore report both unfitted RSA and held-out RDM regression beyond lexical confounds, using item-wise partitions to separate training and test stimuli \citep{schutt2023statistical}. This combination lets us assess which model comparisons recur across evaluations and which are specific to a learned response mapping or a representational score.

\paragraph{Closest work.}
Concurrently with our work, \citet{devarda2026behavioral} compare a low-dimensional effort predictor (surprisal, frequency, and word length) with high-dimensional final-layer embeddings across eight behavioral and five neural datasets, several of which overlap with ours (Natural Stories, Provo, MECO, ZuCo, Tuckute, UCL). They report that behavioral measures are well explained by the effort predictor alone, whereas brain responses require the richer embedding space, and that this split coincides with opposite scaling trends within the GPT-2 family: smaller models fit behavior better, larger models fit the brain better. Their question is largely orthogonal to ours. They hold the model class fixed and ask which \emph{predictor} best explains each modality; we hold the predictors fixed and ask which \emph{model class}, architecturally or developmentally constrained versus standard, best explains each modality, across a much larger sweep of models and datasets. Our size coefficients are consistent with their behavioral--neural scaling contrast for surprisal, and extend it to the question of human-inspired constraints.

Our study thus differs from prior work in two ways. First, rather than taking a side in the constrained-versus-standard debate, we place memory-constrained architectures, BabyLM-trained models, and standard LLMs in the same evaluation pipeline, so that any advantage for the human-inspired models appears as a direct comparison rather than being read across papers with different datasets and metrics. Second, we evaluate this comparison across self-paced reading and eye-tracking corpora together with fMRI, EEG/ERP, MEG, and intracranial recordings, and with both predictive and representational measures, so that conclusions about alignment are not artifacts of any single modality, measure, or evaluation method.
